\documentclass[11pt]{article}

\usepackage[margin=1in]{geometry}
\usepackage{amsmath,amssymb}
\usepackage{booktabs}
\usepackage{enumitem}
\usepackage{hyperref}

\title{Deep Reinforcement Learning Formulation for\\
Orbital Edge Computing Depth Selection}
\author{}
\date{}

\begin{document}

\maketitle

% ============================================================
\section{Problem Overview}
% ============================================================

We consider a deep reinforcement learning (DRL) approach for adaptive
compression-depth selection in an orbital edge computing (OEC) system.
Earth-observation satellites generate image-processing tasks while passing
over predefined areas of interest (AOIs). Each task consists of a collection
of images that must be encoded onboard the source satellite and transmitted
through a time-varying low-Earth-orbit (LEO) satellite network to a ground
base station.

Each task $k \in \mathcal{K}$ is characterized by a source satellite $s_k$,
destination ground station $g_k$, arrival time $a_k$, number of images $N_k$,
importance weight $w_k$, and soft deadline $d_k$. Images are compressed using
a residual-quantized variational autoencoder (RQ-VAE). The compression depth
$q_k$ is selected from
\[
\mathcal{D} = \{2,4,8,16\}.
\]

Increasing the compression depth improves the measured downstream image
quality in the current system but also increases the number of bits that must
be transmitted. The depth-selection problem therefore creates a tradeoff
between image quality and network resource consumption.

The DRL agent learns a task-level policy that balances image quality, timely
delivery, coverage, and resource cost under time-varying network connectivity
and capacity constraints. The implementation uses Proximal Policy
Optimization (PPO) with action masking. The following sections describe the
satellite-network environment, task-generation process, DRL formulation,
unified utility, reward function, PPO training procedure, evaluation
protocol, and relevant implementation details.


% ============================================================
\section{Satellite Network Environment}
% ============================================================

\subsection{LEO Constellation and Simulation Horizon}

The default DRL environment uses the Kuiper-630 LEO constellation. The
constellation contains 34 orbital planes with 34 satellites per plane, for a
total of
\[
|\mathcal{S}| = 34 \times 34 = 1156
\]
satellites. The satellites orbit at an altitude of 630 km with an inclination
of $51.9^\circ$.

The network is simulated for five hours using scheduling slots of duration
\[
\Delta t = 30\text{ s}.
\]

The total configured simulation duration is
\[
T_{\mathrm{sim}} = 18000\text{ s}.
\]

The implementation defines
\[
N_{\mathrm{slots}}
=
\frac{18000}{30}+1
=
601
\]
discrete slot indices, including the initial slot at time zero.


\subsection{Time-Varying Network Topology}

The satellite network is represented by a time-varying graph
\[
\mathcal{G}_t=(\mathcal{V},\mathcal{E}_t),
\]
where the set of feasible communication links $\mathcal{E}_t$ changes as the
satellites move.

Each satellite is assigned inter-satellite links (ISLs) according to a
$+$Grid structure. Satellites connect to their neighboring satellite in the
same orbital plane and to the same-index satellite in the adjacent orbital
plane. Candidate ISLs are enabled or disabled in each slot according to
line-of-sight and maximum-range feasibility constraints.

The configured maximum ISL range is 5016 km. An ISL is also considered
blocked if its line of sight passes below the configured Earth-grazing
clearance.

Ground-to-satellite links (GSLs) are available when a satellite is visible
from a ground station at an elevation angle of at least
\[
20^\circ.
\]

The simulation contains four ground stations:
\begin{itemize}
    \item Tokyo,
    \item New York,
    \item Sao Paulo,
    \item Sydney.
\end{itemize}

For each slot and destination ground station, shortest-path routing is
computed over the currently feasible network using Dijkstra's algorithm.
Physical link distance is used as the graph edge weight. The resulting path
therefore minimizes total propagation distance and, equivalently under a
constant propagation speed, propagation delay.

A task can transmit during a slot only if a feasible route exists between its
source satellite and destination ground station.


\subsection{Network Capacities}

The default DRL experiments use the ground-station-limited scenario. Its
configured rates are
\[
R_{\mathrm{ISL}} = 100\text{ Mbps},
\]
\[
R_{\mathrm{GSL}} = 100\text{ Mbps},
\]
and
\[
R_{\mathrm{GS}} = 2\text{ Mbps}.
\]

Here, $R_{\mathrm{GS}}$ represents the aggregate receiving capacity of a
ground station across satellites communicating with that station.

For an edge or capacity resource $e$ with rate $R_e$, the capacity available
during one scheduling slot is
\[
C_e = R_e\Delta t.
\]

The DRL training procedure can override the ground-station rate to expose the
policy to multiple capacity regimes while retaining the same underlying
satellite topology and task model.


\subsection{Areas of Interest}

The environment contains eight fixed AOIs:
\begin{itemize}
    \item California fires,
    \item Amazon basin,
    \item Sahel,
    \item Ganges delta,
    \item Barrier reef,
    \item Dnipro basin,
    \item Tohoku coast,
    \item Central Andes.
\end{itemize}

An AOI becomes observable when a satellite reaches an elevation angle of at
least
\[
40^\circ
\]
relative to that AOI.


\subsection{Task Generation}

A task is generated when at least one satellite observes an AOI above the
configured imaging-elevation threshold and the AOI's cooldown period has
elapsed.

The cooldown is
\[
1800\text{ s}.
\]

If multiple satellites can observe the AOI during the same slot, the
satellite with the highest elevation angle is selected as the source
satellite.

The destination ground station is chosen as the ground station geographically
nearest to the AOI.

Each task is represented as
\[
k=
(s_k,g_k,a_k,N_k,w_k,d_k),
\]
where:
\begin{itemize}
    \item $s_k$ is the source satellite,
    \item $g_k$ is the destination ground station,
    \item $a_k$ is the arrival slot,
    \item $N_k$ is the number of images,
    \item $w_k$ is the task importance,
    \item $d_k$ is the absolute soft deadline.
\end{itemize}

Under the default ground-station-limited scenario,
\[
N_k
\sim
\mathrm{UniformInteger}(100000,300000).
\]

Task importance is sampled uniformly from
\[
w_k\in\{1,2,3\}.
\]

The deadline offset is sampled between 1800 and 3600 seconds after task
arrival:
\[
d_k
=
a_k\Delta t
+
D_k,
\]
where
\[
D_k
\sim
\mathrm{UniformInteger}(1800,3600).
\]

At most 64 tasks are generated in an episode.


\subsection{Onboard Encoding}

Images become available for transmission gradually as they are encoded
onboard the source satellite.

The measured encoding time per image is
\[
t_{\mathrm{enc}}=12.34\text{ ms}.
\]

Therefore, the encoding rate is approximately
\[
R_{\mathrm{enc}}
=
\frac{1}{0.01234}
\approx
81
\text{ images/s}.
\]

If task $k$ arrived at slot $a_k$, the implementation determines the number
of encoded images available at absolute time $t$ by
\[
N_k^{\mathrm{encoded}}(t)
=
\operatorname{clip}
\left[
(t-a_k\Delta t)R_{\mathrm{enc}},
0,
N_k
\right].
\]


\subsection{RQ-VAE Compression Depth}

The agent selects a compression depth from
\[
\mathcal{D}=\{2,4,8,16\}.
\]

The payload size per image at depth $q$ is
\[
b_q=88q\text{ bytes}.
\]

The unified utility uses measured downstream segmentation quality obtained
from the FLAIR validation set. The quality values are floor-anchored using
the measured downstream segmentation results.

\begin{table}[h]
\centering
\begin{tabular}{ccc}
\toprule
Depth $q$ & Payload per Image & Quality $s_q$ \\
\midrule
2  & 176 B  & 0.182463 \\
4  & 352 B  & 0.262156 \\
8  & 704 B  & 0.312165 \\
16 & 1408 B & 0.360328 \\
\bottomrule
\end{tabular}
\caption{RQ-VAE depth-dependent payload and downstream quality used by the
DRL environment.}
\label{tab:depths}
\end{table}


% ============================================================
\section{Deep Reinforcement Learning Formulation}
% ============================================================

The depth-selection problem is formulated as a sequential decision problem.
The underlying satellite network evolves in fixed 30-second slots, while the
DRL agent operates at task decision points. An agent decision is requested
when an active task has not yet committed to a compression depth.

Once a task commits to a depth, subsequent transmission for that task is
executed automatically by the shared satellite-network simulator.


\subsection{Agent}

The DRL agent controls task-level compression-depth commitment as well as
defer and drop decisions.

When an uncommitted task reaches the front of the decision queue, the agent
receives a compact observation describing the task and its current context.
It then selects one of six discrete actions.

Once a compression depth has been selected for a task, that depth remains
fixed for the remainder of the task.


\subsection{Observation Space}

The agent receives a 16-dimensional continuous observation
\[
\mathbf{o}_t\in\mathbb{R}^{16}.
\]

The observation consists of 12 task-specific features and four
environment-context features.

For current task $k$, the first feature is the fraction of images remaining:
\[
o_1
=
\frac{N_k-r_k}{N_k},
\]
where $r_k$ is the number of images already delivered.

The second feature is the log-scaled task size:
\[
o_2
=
\frac{\log_{10}(\max(N_k,1))}{6}.
\]

The third feature is normalized task importance:
\[
o_3
=
\frac{w_k}{3}.
\]

The fourth feature represents normalized time to deadline:
\[
o_4
=
\operatorname{clip}
\left(
\frac{d_k-t}{D_{\max}},
-2,
2
\right),
\]
where
\[
D_{\max}=3600\text{ s}.
\]

The fifth feature is current task freshness:
\[
o_5=\phi_k(t).
\]

The sixth feature is the fraction of task images that have been encoded:
\[
o_6
=
\operatorname{clip}
\left(
\frac{N_k^{\mathrm{encoded}}(t)}{N_k},
0,
1
\right).
\]

The seventh feature indicates whether a feasible route currently exists:
\[
o_7
=
\begin{cases}
1, & \text{if a route exists},\\
0, & \text{otherwise}.
\end{cases}
\]

The eighth feature is normalized route length:
\[
o_8
=
\begin{cases}
|P_{k,t}|/30, & \text{if a route exists},\\
1, & \text{otherwise},
\end{cases}
\]
where $|P_{k,t}|$ denotes the number of capacity-resource entries associated
with the current route.

The next four features form a one-hot encoding of the committed compression
depth:
\[
(o_9,o_{10},o_{11},o_{12}).
\]

They correspond to depths
\[
(2,4,8,16).
\]

For an uncommitted task, all four depth indicators are zero.

The final four features provide environment-level context.

Normalized simulation time is
\[
o_{13}
=
\frac{t_{\mathrm{slot}}}{N_{\mathrm{slots}}}.
\]

The normalized number of active tasks is
\[
o_{14}
=
\frac{|\mathcal{K}^{\mathrm{active}}_t|}{16}.
\]

The normalized number of active tasks sharing the current task's destination
ground station is
\[
o_{15}
=
\frac{
|\{j\in\mathcal{K}^{\mathrm{active}}_t:g_j=g_k\}|
}{8}.
\]

The final feature is the mean freshness of all currently active tasks:
\[
o_{16}
=
\frac{1}{|\mathcal{K}^{\mathrm{active}}_t|}
\sum_{j\in\mathcal{K}^{\mathrm{active}}_t}
\phi_j(t),
\]
with a value of one when no active tasks are present.

The implementation represents this observation using a Gymnasium
\texttt{Box} space with shape $(16,)$ and component bounds $[-10,10]$.

The observation does not explicitly contain the configured ground-station
bandwidth, individual residual link capacities, or a direct congestion
measurement. Network state is represented indirectly through route
availability, route length, active-task count, destination contention, and
task freshness.


\subsection{Action Space}

The action space is
\[
\mathcal{A}=\{0,1,2,3,4,5\}.
\]

The actions correspond to
\[
\begin{aligned}
0 &\rightarrow q=2,\\
1 &\rightarrow q=4,\\
2 &\rightarrow q=8,\\
3 &\rightarrow q=16,\\
4 &\rightarrow \text{defer},\\
5 &\rightarrow \text{drop}.
\end{aligned}
\]

Actions 0 through 3 commit the current task permanently to the corresponding
compression depth.

The defer action leaves the task uncommitted and removes it from the
remaining decision queue for the current network slot. The task can therefore
return as a decision point in a later slot.

The drop action permanently abandons the task. A dropped task is excluded
from subsequent transmission and scheduling decisions.


\subsection{Action Masking}

Action masking prevents the policy from committing a task to a compression
depth when no feasible route currently exists.

Let
\[
m_t(a)\in\{0,1\}
\]
indicate whether action $a$ is valid.

When a feasible route exists,
\[
m_t(a)=1
\qquad
\forall a\in\mathcal{A}.
\]

When no route exists,
\[
m_t(a)
=
\begin{cases}
1, & a\in\{4,5\},\\
0, & a\in\{0,1,2,3\}.
\end{cases}
\]

Thus, defer and drop are always available, while depth commitment requires
current route feasibility.


\subsection{Transition Process}

At each network slot, the environment constructs the set of active tasks:
\[
\mathcal{K}^{\mathrm{active}}_t
=
\left\{
k:
a_k\le t,\;
r_k<N_k,\;
k\text{ not dropped}
\right\}.
\]

Active tasks are ordered by deadline.

If the next task has no committed compression depth, execution pauses and the
DRL agent is asked to select an action.

Once a depth is committed, transmission is performed using the shared network
execution logic also used by the non-RL schedulers.

For task $k$, the simulator first obtains the current shortest path
\[
P_{k,t}
\]
from source satellite $s_k$ to destination ground station $g_k$.

The path is converted into capacity resources. These include each traversed
ISL, the final satellite-to-ground link, and the aggregate receiving capacity
of the destination ground station.

For resource $e$, slot capacity is
\[
C_e=R_e\Delta t.
\]

At depth $q_k$, one image contains
\[
B_{q_k}=8b_{q_k}
\]
bits.

The number of images that can be delivered during a slot is limited by both
the number of encoded but undelivered images and the bottleneck residual
capacity along the path:
\[
y_{k,t}
=
\min
\left\{
N_k^{\mathrm{encoded}}(t)-r_k,\;
\min_{e\in P_{k,t}}
\frac{C^{\mathrm{res}}_{e,t}}{B_{q_k}}
\right\}.
\]

The corresponding number of bits is deducted from the residual capacity of
every resource along the path.

After transmission,
\[
r_k
\leftarrow
r_k+y_{k,t}.
\]

The delivery arrival time is taken to be the end of the current scheduling
slot plus propagation delay:
\[
t_{\mathrm{arrive}}
=
(t_{\mathrm{slot}}+1)\Delta t
+
d^{\mathrm{prop}}_{k,t}.
\]

The delivery's utility is evaluated at this arrival time.

Tasks that have already committed to a depth may execute automatically while
the simulator advances. When the current slot's decision queue is exhausted,
the environment advances to the next slot and continues until another
uncommitted task requires an agent decision or the episode ends.


\subsection{Episode Definition}

One episode spans the complete configured satellite-network simulation.

The environment terminates when the slot index reaches the configured
simulation horizon. Under the default configuration, the underlying network
scenario covers five hours.

The environment does not use a separate truncation condition.


% ============================================================
\section{Unified Utility and Reward Formulation}
% ============================================================

\subsection{Task Utility}

The unified utility defines the realized value of task $k$ as
\[
U_k
=
\bar{w}_k
\left[
\omega_Qs_{q_k}\widehat{G}_k
-
\omega_TT_k
-
\omega_EC_k
\right].
\]

Here,
\begin{itemize}
    \item $\bar{w}_k$ is normalized task importance,
    \item $s_{q_k}$ is downstream quality at compression depth $q_k$,
    \item $\widehat{G}_k$ is freshness-weighted concave coverage,
    \item $T_k$ is tardiness,
    \item $C_k$ is normalized resource cost.
\end{itemize}

In unified mode,
\[
\bar{w}_k
=
\frac{w_k}{\max(\{1,2,3\})}
=
\frac{w_k}{3}.
\]

The utility weights are
\[
\omega_Q=1.0,
\qquad
\omega_T=0.25,
\qquad
\omega_E=0.05,
\qquad
\omega_F=0.10.
\]


\subsection{Downstream Quality}

The unified objective uses downstream semantic-segmentation quality.

Let $m_q$ denote the measured mIoU at compression depth $q$,
$m_{\mathrm{ref}}$ the uncompressed/reference segmentation mIoU, and
$m_{\mathrm{floor}}$ the segmentation performance when the relevant image
information is unavailable.

The floor-anchored quality score is
\[
s_q
=
\frac{m_q-m_{\mathrm{floor}}}
{m_{\mathrm{ref}}-m_{\mathrm{floor}}}.
\]

The measured values used by the current four depth actions are
\[
s_2=0.182463,
\qquad
s_4=0.262156,
\qquad
s_8=0.312165,
\qquad
s_{16}=0.360328.
\]


\subsection{Freshness}

Task freshness remains one through the soft deadline and then decays
exponentially:
\[
\phi_k(t)
=
\begin{cases}
1,
&
t\le d_k,\\[4pt]
\exp[-\alpha(t-d_k)],
&
t>d_k.
\end{cases}
\]

The configured decay rate is
\[
\alpha=\frac{1}{300}\text{ s}^{-1}.
\]

Thus, an image delivered on time receives full freshness credit, while the
value of increasingly late deliveries decays exponentially.


\subsection{Concave Coverage}

Let
\[
\rho_k=\frac{r_k}{N_k}
\]
denote the delivered fraction of task $k$.

Unified utility uses a concave coverage function so that the first portion of
a task delivered has greater marginal value than the final portion.

The implementation uses four piecewise-linear segments:
\[
(\Delta\rho_j,m_j)
\in
\{
(0.25,2.00),
(0.25,1.12),
(0.25,0.60),
(0.25,0.28)
\}.
\]

Each segment covers one quarter of the task, while its slope decreases as
coverage increases.

For a delivery covering task fraction $\Delta\rho$ at freshness
$\phi_k(t)$, the incremental coverage gain within segment $j$ is
\[
\Delta\widehat{G}_k
=
m_j\phi_k(t)\Delta\rho.
\]

The implementation fills these segments chronologically. Consequently,
early coverage of a previously unserved task receives greater marginal
utility than the final portion of an already well-covered task.


\subsection{Tardiness}

For an image arriving at time $t$, the normalized tardiness coefficient is
\[
\tau_k(t)
=
\min
\left(
1,
\frac{\max(0,t-d_k)}
{T_{\mathrm{ref}}}
\right),
\]
where
\[
T_{\mathrm{ref}}=300\text{ s}.
\]

A delivery at or before the deadline therefore has zero tardiness penalty.
After the deadline, the coefficient grows linearly until reaching one.


\subsection{Resource Cost}

The resource-cost coefficient for compression depth $q$ is
\[
c_q
=
\omega_{\mathrm{tx}}
\frac{b_q}{b_{\max}}
+
\omega_{\mathrm{enc}},
\]
where
\[
\omega_{\mathrm{tx}}=0.5,
\qquad
\omega_{\mathrm{enc}}=0.5.
\]

The first component increases with transmitted payload size. The second
represents the fixed per-image encoding component.

Because the measured encoding time is effectively constant across the
available compression depths, the encoding component does not distinguish
between depth choices, whereas the transmission component increases with
depth.


\subsection{Incremental Realized Utility}

Suppose a delivery contains a fraction
\[
\Delta\rho=\frac{y_{k,t}}{N_k}
\]
of task $k$.

Its incremental realized task utility is
\[
\Delta U_k
=
\bar{w}_k
\left[
\omega_Qs_{q_k}\Delta\widehat{G}_k
-
\omega_T\tau_k(t)\Delta\rho
-
\omega_Ec_{q_k}\Delta\rho
\right].
\]

The task's delivered utility accumulates these contributions over all of its
deliveries.


\subsection{Fairness}

After task utilities have been accumulated, the final unified system utility
is
\[
U
=
\sum_{k\in\mathcal{K}}U_k
+
\omega_F|\mathcal{K}|
\min_{k\in\mathcal{K}}\widehat{U}_k.
\]

The normalized task utility used for fairness is
\[
\widehat{U}_k
=
\frac{U_k}
{\bar{w}_k\omega_Qs_{\max}},
\]
where
\[
s_{\max}
=
\max_{q\in\mathcal{D}}s_q.
\]

This normalization makes the fairness floor less dependent on the absolute
scale and importance weight of a task.

The minimum normalized task utility
\[
\min_k\widehat{U}_k
\]
acts as a maximin fairness surrogate. It discourages a scheduler from
obtaining high aggregate utility while completely starving an individual
task.


\subsection{Jain Fairness Diagnostic}

The implementation also reports Jain's fairness index:
\[
J
=
\frac{
\left(\sum_k\widehat{U}_k\right)^2
}{
|\mathcal{K}|
\sum_k\widehat{U}_k^2
}.
\]

Jain's index is reported as a diagnostic metric only. It is not directly
optimized by the unified objective.


\subsection{Backlog Potential}

In addition to realized task utility, PPO receives a potential-based reward
shaping term based on the remaining network backlog.

Let $B_t$ denote the estimated remaining number of bits across active,
non-dropped tasks.

For a task with committed depth $q_k$, its remaining backlog is estimated
using that depth's payload size. For an uncommitted task, the implementation
uses depth 16 when estimating backlog.

The potential is
\[
\Phi_t
=
-10^{-3}
\frac{B_t}{10^{11}}.
\]

A smaller backlog makes this negative potential closer to zero.


\subsection{PPO Reward}

Let
\[
\Delta J_t
=
U_k^{\mathrm{after}}
-
U_k^{\mathrm{before}}
\]
represent the immediate change in the current task's delivered utility
associated with the current decision and execution.

The PPO reward is
\[
r_t
=
\Delta J_t
+
\gamma_{\Phi}\Phi_{t+1}
-
\Phi_t,
\]
where
\[
\gamma_{\Phi}=0.999.
\]

The backlog potential provides an additional learning signal that encourages
reduction of outstanding network work.

An important distinction is that the final schedule-level maximin fairness
bonus is computed when the complete collection of tasks is evaluated. It is
not directly added to the per-step PPO reward.

Thus, PPO directly receives realized quality, freshness/coverage, tardiness,
resource-cost, and backlog-shaping signals, while the global fairness term is
applied to the completed schedule during unified-utility evaluation.


% ============================================================
\section{PPO Implementation and Training}
% ============================================================

\subsection{MaskablePPO}

The policy is trained using MaskablePPO from the \texttt{sb3-contrib}
extension to Stable-Baselines3.

MaskablePPO combines PPO with state-dependent action masks. This ensures that
actions identified as invalid by the environment are excluded from the
policy's available choices.

PPO optimizes a clipped policy objective. Let
\[
r_t(\theta)
=
\frac{
\pi_\theta(a_t|\mathbf{o}_t)
}{
\pi_{\theta_{\mathrm{old}}}(a_t|\mathbf{o}_t)
}
\]
denote the probability ratio between the updated and previous policies.

Given estimated advantage $\widehat{A}_t$, the standard clipped PPO surrogate
objective is
\[
L^{\mathrm{CLIP}}(\theta)
=
\mathbb{E}_t
\left[
\min
\left(
r_t(\theta)\widehat{A}_t,
\operatorname{clip}
\left(
r_t(\theta),
1-\epsilon,
1+\epsilon
\right)
\widehat{A}_t
\right)
\right].
\]

The action mask further restricts the policy distribution to actions that
are valid at the current decision point.


\subsection{Policy Network}

The implementation uses an MLP policy.

The configured hidden-layer architecture is
\[
[128,128].
\]

Conceptually, the observation is processed as
\[
16
\rightarrow
128
\rightarrow
128
\rightarrow
\text{policy/value outputs}.
\]


\subsection{Training Hyperparameters}

The principal PPO hyperparameters are shown in
Table~\ref{tab:ppo-hyperparameters}.

\begin{table}[h]
\centering
\begin{tabular}{lc}
\toprule
Parameter & Value \\
\midrule
Algorithm & MaskablePPO \\
Policy & MLPPolicy \\
Observation dimension & 16 \\
Number of actions & 6 \\
Hidden architecture & $[128,128]$ \\
Rollout length per environment & 2048 \\
Batch size & 256 \\
Learning rate & $3\times10^{-4}$ \\
Discount factor $\gamma$ & 0.999 \\
GAE parameter $\lambda$ & 0.95 \\
Entropy coefficient & 0.01 \\
Nominal training timesteps & 5,000,000 \\
\bottomrule
\end{tabular}
\caption{MaskablePPO training configuration.}
\label{tab:ppo-hyperparameters}
\end{table}


\subsection{Training Distribution}

Training exposes the policy to three ground-station capacity settings:
\[
R_{\mathrm{GS}}
\in
\{1.5,2.0,2.5\}
\text{ Mbps}.
\]

For each capacity setting, task-generation seeds
\[
0,1,\ldots,99
\]
are used.

This produces
\[
3\times100=300
\]
training environments.

The seed controls the deterministic pseudo-random generation of task
properties such as image counts, task importance, and deadline offsets.

The environments are combined using a vectorized environment.

Each environment contributes
\[
n_{\mathrm{steps}}=2048
\]
transitions to a rollout.

Therefore, a complete vectorized rollout contains
\[
300\times2048
=
614400
\]
transitions.

Because PPO completes full rollouts, the actual number of training
transitions can exceed the nominal requested training budget of five million
timesteps.


\subsection{Per-Environment Capacity Configuration}

Each training environment retains its assigned ground-station rate during
environment construction, reset, action-mask evaluation, and stepping.

This is necessary because the shared capacity functions read the current
configuration when transmission is executed. Maintaining the override for
the complete environment interaction ensures that the intended
1.5, 2.0, and 2.5 Mbps training regimes are actually represented during
training.


\subsection{Legacy and Unified PPO}

The same DRL architecture can be trained using either the legacy or unified
utility configuration.

The observation space, action space, action masking, transition process,
policy architecture, and PPO hyperparameters remain unchanged.

In legacy mode,
\[
\omega_Q=1,
\qquad
\omega_T=\omega_E=\omega_F=0.
\]

Legacy mode uses one linear coverage segment,
\[
(1.0,1.0),
\]
does not normalize task weights, and uses reconstruction quality
\[
s_q=1-\mathrm{LPIPS}_q.
\]

The legacy reconstruction-quality values for depths
$q\in\{2,4,8,16\}$ are
\[
0.4430,\quad
0.4686,\quad
0.4829,\quad
0.4961.
\]

Unified mode instead enables normalized task weights, downstream segmentation
quality, concave coverage, tardiness, resource cost, and the final maximin
fairness term.

This setup supports a controlled comparison between a PPO policy trained
using the historical objective and a PPO policy trained using the unified
objective.


% ============================================================
\section{Evaluation Protocol}
% ============================================================

\subsection{Held-Out Task Seeds}

Training uses task-generation seeds 0 through 99.

Evaluation uses held-out seeds
\[
100,101,\ldots,119.
\]

These seeds generate task realizations that are not included in PPO training.


\subsection{In-Distribution Evaluation}

The nominal in-distribution evaluation uses
\[
R_{\mathrm{GS}}=2.0\text{ Mbps}.
\]

This capacity lies within the set of rates presented during training.


\subsection{Bandwidth Generalization}

The policy can additionally be evaluated across
\[
R_{\mathrm{GS}}
\in
\{0.75,1.0,2.0,3.0,4.0\}
\text{ Mbps}.
\]

The 0.75, 1.0, 3.0, and 4.0 Mbps conditions are outside the exact set of
training rates and can therefore be used to examine generalization to
different bandwidth regimes.


\subsection{Baseline Comparison}

PPO policies are evaluated using the same satellite-network simulation and
utility calculation used by the non-RL schedulers.

This permits matched comparisons with implemented scheduling strategies such
as fixed-depth policies, greedy schedulers, and model-predictive-control
methods.

Relevant evaluation quantities include:
\begin{itemize}
    \item total unified utility,
    \item delivered fraction,
    \item quality contribution,
    \item tardiness contribution,
    \item resource-cost contribution,
    \item minimum normalized task utility,
    \item Jain's fairness index.
\end{itemize}


% ============================================================
\section{Implementation Details and Current Limitations}
% ============================================================

\subsection{Shared Network Execution}

The DRL environment does not use a separate network model from the other
schedulers.

After PPO commits a task to a compression depth, the task uses the shared
route construction, edge-capacity accounting, propagation-delay calculation,
and delivery-recording mechanisms.

This provides consistent execution semantics between the learned policy and
the non-RL scheduling baselines.

A route consumes capacity from each traversed ISL, the final GSL, and the
destination ground station's aggregate receive budget.


\subsection{Decision-Point Time Scale}

An RL step does not necessarily correspond to one 30-second network slot.

The environment may execute already committed tasks and advance across
network slots internally before reaching another task requiring an agent
decision.

The agent therefore operates at task decision points layered on top of the
fixed-slot network simulation.


\subsection{Training Reward Versus Final Utility}

The PPO training reward and final unified schedule utility are related but
are not identical.

PPO receives the current task's immediate realized utility change together
with backlog-based potential shaping.

The final schedule-level evaluation additionally includes the global maximin
fairness term
\[
\omega_F|\mathcal{K}|
\min_k\widehat{U}_k.
\]

Therefore, fairness affects the final unified score but is not directly
provided as a per-step PPO reward.


\subsection{Temporal Credit Assignment}

After an agent action, the environment may automatically execute tasks whose
depths were committed previously while advancing toward the next decision
point.

The immediate reward term
\[
\Delta J_t
\]
is computed from the current task's delivered utility around the explicit
agent action. Utility produced by subsequent automatic transmissions is
reflected in task state and final schedule evaluation, but it is not
necessarily attributed directly to that action's immediate $\Delta J_t$.

This is an implementation detail of the current temporal credit-assignment
scheme.


\subsection{Bandwidth Observability}

The current observation does not explicitly contain the configured
ground-station capacity
\[
R_{\mathrm{GS}},
\]
nor does it directly expose residual capacities for individual links.

Consequently, two decision points with similar task and route features can
produce similar observations even when they occur under substantially
different bandwidth regimes.

The policy receives indirect indicators of network conditions through route
availability, route length, active-task count, destination contention, and
freshness. These features do not, however, uniquely determine the available
network capacity.

Explicit bandwidth, residual-capacity, or congestion features represent a
possible extension for improving adaptation across substantially different
network-capacity regimes.


% ============================================================
\section{Summary}
% ============================================================

The implemented DRL system separates learned task-level decision making from
the physical satellite-network simulator.

The interaction can be summarized as
\[
\boxed{
\begin{array}{c}
\text{Satellite motion and AOI visibility generate tasks}\\
\downarrow\\
\text{Environment constructs a 16-dimensional observation}\\
\downarrow\\
\text{Invalid depth actions are masked}\\
\downarrow\\
\text{PPO selects depth, defer, or drop}\\
\downarrow\\
\text{Network simulator executes feasible transmission}\\
\downarrow\\
\text{Delivered images update task utility}\\
\downarrow\\
\text{Incremental utility and backlog shaping form the reward}\\
\downarrow\\
\text{Environment advances to the next task decision point}
\end{array}
}
\]

The learned policy determines compression-depth and admission decisions,
while the satellite-network simulator determines orbital connectivity,
routing, capacity-constrained transmission, encoding availability,
propagation delay, and realized delivery outcomes.

The formulation therefore allows PPO to learn adaptive depth-selection
behavior while preserving the same physical network execution model used by
the existing scheduling baselines.

\end{document}
